\model{List and Pair}
  \begin{center}
    \small
    \begin{tabular}{p{3.0in} p{2.8in}}
      \begin{minipage}{3.0in}
        \begin{cpplst}
#include <list>
#include <utility>

list<string> tasks =
    {"breakfast", "class", "lunch"};

tasks.push_front("wake up");
tasks.push_back("dinner");
tasks.pop_front();

for (string task : tasks) {
    cout << task << endl;
}
        \end{cpplst}
      \end{minipage}
      &
      \begin{minipage}{2.8in}
        \vskip 10pt
        {\bf After push\_front("wake up"):}\\[4pt]
        \begin{tabular}{|c|c|c|c|c|}
          \hline
          wake up & breakfast & class & lunch \\
          \hline
        \end{tabular}\\[8pt]
        {\bf After push\_back("dinner"):}\\[4pt]
        \begin{tabular}{|c|c|c|c|c|}
          \hline
          wake up & breakfast & class & lunch & dinner \\
          \hline
        \end{tabular}\\[8pt]
        {\bf After pop\_front():}\\[4pt]
        \begin{tabular}{|c|c|c|c|}
          \hline
          breakfast & class & lunch & dinner \\
          \hline
        \end{tabular}\\[10pt]
        {\bf Output:}
        \hrule\vskip 5pt
        {\tt breakfast}\\
        {\tt class}\\
        {\tt lunch}\\
        {\tt dinner}
      \end{minipage}
    \end{tabular}
    \vskip 12pt
    \renewcommand{\arraystretch}{1.3}
    \begin{tabular}{|c|c|c|}
      \hline
      \rowcolor{orange!20}
      \textbf{Operation} & \textbf{vector} & \textbf{list} \\
      \hline
      \cpp{push_back(val)}  & yes & yes \\
      \hline
      \cpp{pop_back()}      & yes & yes \\
      \hline
      \cpp{push_front(val)} & no  & yes \\
      \hline
      \cpp{pop_front()}     & no  & yes \\
      \hline
      \cpp{at(i)}           & yes & no  \\
      \hline
      range-based for       & yes & yes \\
      \hline
    \end{tabular}
    \vskip 12pt
    \begin{minipage}{3.5in}
      \begin{cpplst}
pair<string, int> student = {"Alice", 95};
cout << student.first << " scored ";
cout << student.second << endl;
      \end{cpplst}
    \end{minipage}
    \quad
    \begin{minipage}{2in}
      {\bf Output:}
      \hrule\vskip 5pt
      {\tt Alice scored 95}
    \end{minipage}
  \end{center}

  {\it\large Refer to Model 2 above as your team develops consensus answers
    to the questions below.}

  \quest{15 min}

  \Q What {\tt \#include} header is required to use each of the following?
    \begin{enumerate}
      \itemsep 10pt
      \item \cpp{list} \hfill\ans[3.5in]{\cpp{#include <list>}}
      \item \cpp{pair} \hfill\ans[3.5in]{\cpp{#include <utility>}}
    \end{enumerate}

  \vskip -20pt

  \Q Using the operations table in the model, answer the following.
    \begin{enumerate}
      \itemsep 10pt
      \item Write a line of C++ to access the 3rd element (index~2)
        of a {\tt vector<int>} named {\tt v}.
        \hfill\ans[1in]{\cpp{v.at(2)}}

      \item Can you do the same operation with a {\tt list}?
        Explain why or why not.
        \begin{answer}[0.5in]
          No. A \cpp{list} does not support \cpp{at(i)} because it is
          a linked structure; elements are not stored contiguously,
          so random access by index is not supported.
        \end{answer}

      \item A {\tt vector} does not support {\tt push\_front}.  Why
        might adding an element to the front of a vector be slow compared
        to a list?
        \begin{answer}[0.5in]
          Every existing element must be shifted one position to the
          right to make room at index~0, which takes time proportional
          to the size of the vector.  A list can insert at the front
          in constant time by adjusting pointers.
        \end{answer}
    \end{enumerate}

  \vskip -20pt

  \Q When would you choose a {\tt list} over a {\tt vector}?  \key\\[-2mm]
    When would you choose a {\tt vector} over a {\tt list}?
    \begin{answer}[0.75in]
      Use a \cpp{list} when you need to frequently add or remove
      elements from both ends, or when random access is not needed.
      Use a \cpp{vector} when you need fast random access by index
      or when you mostly add/remove from the back.
    \end{answer}


  \Q A {\tt pair} stores two related values of potentially different
    types and is accessed via {\tt .first} and {\tt .second}.
    \begin{enumerate}
      \itemsep 10pt
      \item In the pair declaration above, what is the type of
        {\tt student.first}?
        \hfill\ans[1.5in]{\cpp{string}}

      \item What is the type of {\tt student.second}?
        \hfill\ans[1.5in]{\cpp{int}}

      \item Write one line of C++ to declare a {\tt pair} that stores
        a city name (string) and its population (int), with the city
        {\tt "Bellingham"} and population {\tt 90000}.
        \begin{answer}[0.5in]
          \cpp{pair<string, int> city = {"Bellingham", 90000};}
        \end{answer}

      \item The file {\tt activity28b.cpp} contains the list and pair
        examples.  Run it and verify the output matches the model.
        What does {\tt tasks.front()} return?
        \hfill\ans[1.5in]{\cpp{"breakfast"}}
    \end{enumerate}
